\begin{annexitem}{annex:p7.mr.1}{thm:p07-system-t-numerical-computability}{Numerical values}
\begin{proof}
Use the compatible beta and recursor reductions of the fixed term
calculus. For closed typed terms define computability by induction on
the type. At ground type, a term is computable when there exists a
finite reduction to a numeral. At \(\sigma\to_{\mathrm{ty}}\tau\),
it is computable when its application to every closed computable term
of type \(\sigma\) is computable at type \(\tau\).

This property is closed under inverse reduction: if \(t\to t'\)
and \(t'\) is computable, then \(t\) is computable. At ground type
prepend that step to a finite path. At an arrow type, for any computable
\(u\), compatibility gives \(tu\to t'u\); use the induction
hypothesis at the result type. The same conclusion holds for a finite
sequence of reductions by iteration.

Zero is computable. Successor maps a reduction to a numeral to a
reduction to its successor numeral, so it is computable at its arrow
type. The recursor at \(\sigma\) is computable as well. Take
computable \(a,b,n\) of its three argument types. A finite reduction
of \(n\) to \(\overline k\) gives a reduction of
\(R_\sigma a b n\) to \(R_\sigma a b\overline k\).
Induct on \(k\). At zero the latter reduces to \(a\). At a
successor it reduces to
\(b\overline k(R_\sigma a b\overline k)\), computable by the
induction hypothesis and the computability of \(b\). Inverse
closure proves the required result at every \(k\).

The usual substitution induction on a typed term now says that
replacing its free variables by closed computable terms gives a
computable term. The variable and application cases are immediate;
constants were just checked. For an abstraction, apply it to an
arbitrary computable argument and use its beta reduction and the
induction hypothesis for the body with that additional substitution.
Consequently every closed ground term reduces to a numeral.

The compatible one-step reduction relation is effectively enumerable,
with finitely many candidate redex positions in each finite term.
Breadth-first enumeration of its finite reduction paths therefore
finds a numeral whenever one is reachable. Equivalently, finite syntax
and single reductions have primitive-recursive code checks, so search
through coded finite reduction certificates is a partial-recursive
procedure. \Cref{thm:p04-tm-partial-recursive} realizes that procedure
in the book's fixed machine model. The preceding argument
ensures this search terminates on a well-typed closed ground term.

For uniqueness of its value we need not assume a separate confluence
theorem. Interpret the finite types by \(\mathbb N\) and full
function sets, lambda abstraction and application by their ordinary
set-theoretic operations, and each recursor by recursion on the
natural numbers. Substitution proves the beta equation sound; natural
recursion proves both recursor equations sound. Compatible reductions
therefore preserve denotation. Distinct numerals denote distinct
natural numbers, so any two numerals found by the search agree.
Applying this algorithm to \(f\overline n\), for closed
\(f:0\to_{\mathrm{ty}}0\), computes a total numerical function.
The numerical successor and equality axioms of the quantifier-free
target also hold in this interpretation, as does open induction by
ordinary induction on \(\mathbb N\). Hence its proved term equations,
including the tuple-recursion equations, hold in this same model.
\end{proof}
\end{annexitem}

\begin{annexitem}{annex:p7.mr.2}{thm:p07-modified-realizability-soundness}{Typed witnesses}
\begin{proof}
The positive computational operations in the ordinary realizability
proof have typed total counterparts. We specify them here to distinguish
this interpretation from both program indices and proof-carrying data.
All tuples are lists of terms; a missing input tuple means no lambda
arguments and a missing output tuple means no components. Use the closed
default term at each finite type for unused components. Numerical
conditionals are those of \cref{thm:p07-t-matrix-decision}.

\emph{Application and composition.} Modus ponens applies the extracted
functional for \(A\to B\) to the extracted \(A\)-tuple. Its
modified-realizability clause gives the \(B\)-condition. Composition
of \(A\to B\) and \(B\to C\) is componentwise composition of
those functionals. This remains valid when an intermediate tuple is
empty: the corresponding condition is a proposition with no witness
arguments, rather than an omitted premise.

\emph{Contractions and structural operations.} To realize
\(A\to A\land A\), duplicate its input tuple. To realize
\(A\lor A\to A\), use the numerical tag to select one of its
two input tuples by typed conditionals. A realizing input has tag zero
or one and the selected tuple realizes \(A\). Disjunction injection
uses tag zero, the input tuple, and defaults for the unused branch;
conjunction projection selects its relevant tuple. The commutation
rules exchange tuples and, for disjunction, exchange tags zero and one.
For disjunction monotonicity, leave the retained branch unchanged and
apply the given \(A\to B\) functional to the other branch. These
constructions verify groups 3--6 of the logical presentation without a
truth test of an arbitrary arithmetic formula.

\emph{Implication and falsity.} Currying and uncurrying regroup the
successive tuple arguments of a lambda term; the defining universal
conditions for implication are identical after this regrouping. An
empty output tuple requires no function components but retains the
universal semantic implication. Falsity has no realizing tuple, so
any default output realizes falsity elimination.

\emph{Quantifiers.} From a realizer \(F(z,\vec a)\) for
\(A\to B(z)\), with \(z\) not free in \(A\), abstract \(z\)
to obtain a realizer for \(A\to\forall zB(z)\). Its defining
condition holds for every number by the child theorem's uniformity.
Universal elimination evaluates the input functions at the value of
the specified numerical term. Existential introduction pairs that
value with the input tuple. From a realizer for \(A(z)\to B\),
where \(z\) is not free in \(B\), a realizer for
\(\exists zA(z)\to B\) takes the supplied witness number and
its tuple and applies the former realizer. These are respectively
logical groups 9--12 and have exactly their original eigenvariable
conditions. Together with the preceding operations they cover all
twelve logical groups.

\emph{Equality and numerical axioms.} An atomic formula has no
positive data; its modified-realizability condition is its numerical
truth. Equality replacement on numerical atoms and the successor,
addition and multiplication axioms therefore have the required
conditions in the standard numerical structure. For predecessor
existence use the numerical predecessor as the existential witness.
Its condition is required only when the number is nonzero, where its
successor is that number. Universal closures are handled by the
quantifier constructions.

\emph{Induction.} A realizing input for the induction antecedent
consists of a tuple \(\vec a\) realizing \(A(0)\) and a tuple
\(\vec F\) satisfying the universal step condition. Define
\[
 \vec X(0)=\vec a,\qquad
 \vec X(Si)=\vec F(i,\vec X(i))
\]
by \cref{thm:p07-t-simultaneous-recursion}. The tuple types depend
on the fixed shape of \(A\), not on the value of \(i\). Its terms are in the
arrow-only System T, and their equations hold in the full finite-type
interpretation by the preceding lemma. Induction on the actual natural
number \(i\) proves \(\operatorname{mr}_{A(i)}(\vec X(i))\).
Thus \(\lambda i.\vec X(i)\) realizes the conclusion of the
induction axiom. This also covers an empty tuple of components, in
which case the step condition itself implies the next proposition.

Induction on the finite HA proof now builds the specified terms and
verifies their semantic conditions. Retain any extra numerical
parameters of intermediate premises during the eigenvariable operations;
afterwards set remaining parameters absent from the conclusion to zero.
This preserves the uniform condition for that conclusion. All positive
witnesses of a closed conclusion are therefore closed terms. Only
lambda abstraction, application, default terms, numerical conditionals,
and the proved tuple recursors have been used.
\end{proof}
\end{annexitem}

\begin{annexitem}{annex:p7.mr.3}{thm:p07-bounded-modified-realizability}{Bounded witnesses}
\begin{proof}
Proceed simultaneously through the bounded formula, constructing its
numerical truth test and a canonical tuple, and proving both directions
of the stated truth comparison. Numerical atoms are tested by
\cref{thm:p07-t-matrix-decision}; their tuple is empty. Conjunction
concatenates canonical tuples and combines truth tests. For disjunction,
test its two constituents, choose a true branch when one exists, and
supply the canonical tuple of that branch and defaults for the unused
one. A realizing tuple always identifies a true constituent by the
induction hypothesis.

For implication \(C\to D\), use the constant function with value
the canonical tuple for \(D\). This is valid whenever the implication
is true: if \(C\) is false there is no realizing input by the
induction hypothesis, while if \(D\) is true its canonical tuple
realizes it. Conversely, if \(C\) is true and \(D\) false, the
canonical \(C\)-tuple is a realizing input but there is no possible
realizing output for \(D\). Thus no functional realizes that false
implication. The Boolean truth test is the ordinary implication of
the two already constructed tests. Negation is the special case with
false conclusion. These arguments also handle empty tuple types.

For explicit bounded guards one may write \(x<t\) as
\(\exists u\,(x+Su=t)\) in the arithmetic language. It is
realizable exactly when it is numerically true, and in that case
\(p(t\mathbin{\dot{-}}x)\) is a canonical numerical witness.
Otherwise no number can make its atomic matrix true. Truncated
subtraction and predecessor are System-T terms. The analogous
nonstrict guard has witness \(t\mathbin{\dot{-}}x\).

A bounded universal formula is read as
\(\forall x(x<t\to C(x))\), with \(x\) absent from the bound.
Use the functions which, at \(x\) and a guard witness, return the
canonical tuple for \(C(x)\), ignoring the guard witness. These
are total functions at every \(x\). Below the bound the tuple is
valid whenever the bounded universal is true; outside the bound its
implication condition is vacuous. Conversely, a false bounded universal
has a counterexample below its bound. The guard has a canonical witness
there, and any proposed universal realizer would have to supply a
realizer of the false constituent, which is impossible.

For \(\exists x<t\,C(x)\), use a bounded numerical scan of the
already constructed truth test to find the first satisfying \(x\),
if there is one. Return that number, its canonical guard witness, and
the canonical \(C(x)\)-tuple. If none exists return default data;
no validity is claimed in that case. Every realizing existential tuple,
on the other hand, supplies an actual number below the bound and a true
constituent. Bounded universal and existential truth tests are finite
Boolean folds over the same numerical bound. They are defined by
primitive recursion, and the canonical tuples use only typed
conditionals, bounded scans, and lambda abstraction. Thus all their
components are System-T terms, total even at parameters where the
formula is false.

Finally at each \(t\), the truth test for \(B(\vec n,t)\)
chooses tag zero with its canonical tuple if it is true, and tag one
otherwise. In the latter case no tuple realizes \(B\), so the empty
negative tuple realizes \(\neg B\). Abstracting \(t\) gives the
asserted uniform witness for the universal decidability statement.
\end{proof}
\end{annexitem}

\begin{annexitem}{annex:p7.mr.4}{thm:p07-markov-realizability-boundary}{The halting obstruction}
\begin{proof}
Let \(K\) be the computably enumerable, undecidable self-application
domain set of \cref{thm:halting-undecidable}. Use the bounded-computation
presentation in the proof of \cref{thm:p04-ce-sigma-one}. A fixed
index for a partial computable function with domain \(K\) gives a
bounded arithmetic formula \(H(e,t)\) such that
\[
 e\in K\quad\longleftrightarrow\quad
 \mathcal N\models\exists t\,H(e,t).
 \tag{H}
\]
Concretely this is the displayed bounded computation verifier with
that fixed program numeral in its index argument, \(e\) in its
input argument, and \(t\) its computation certificate. Thus the
matrix really is bounded in the handbook's marked bounded-quantifier
convention; an arbitrary primitive-recursive predicate has not simply
been declared bounded.

Abbreviate
\[
 E(e)=\exists t\,H(e,t),\qquad
 D(e)=\forall t\,(H(e,t)\lor\neg H(e,t)).
\]
By \cref{thm:p07-bounded-modified-realizability}, there is a tuple
of System-T terms \(\vec d(e)\) realizing \(D(e)\), uniformly in
\(e\). Suppose HA proved the universal closure of the corresponding
Markov instance,
\[
 \forall e\bigl(D(e)\to(\neg\neg E(e)\to E(e))\bigr).
 \tag{M}
\]
Modified-realizability soundness supplies a closed tuple of System-T
terms for (M). Apply its \(e\)-component to \(\vec d(e)\).
The positive tuple for \(\neg\neg E(e)\) is empty, so no further
positive argument is required. The resulting tuple has the witness
types of \(E(e)\); in particular its first component is a numerical
System-T term \(f(e)\), defined for every \(e\).

If \(e\in K\), some \(H(e,t)\) is true. The bounded-formula
lemma supplies a modified realizer of that instance and hence one of
\(E(e)\). Its double negation therefore has its true, empty-tuple
modified-realizability condition. More explicitly, negation says that
no tuple realizes its argument, so the double negation says that it is
not the case that no tuple realizes \(E(e)\); this holds for the
witness just supplied. The implication condition for the extracted
realizer of (M) now forces the output tuple to realize \(E(e)\).
Its first component is \(f(e)\), and bounded truth reflection gives
\(H(e,f(e))\). Conversely \(H(e,f(e))\) always implies \(e\in K\)
by (H). We have obtained
\[
 e\in K\quad\longleftrightarrow\quad H(e,f(e)).
\]
By \cref{thm:p07-system-t-numerical-computability}, the closed
abstraction \(\lambda e.f(e)\) computes a total recursive function.
The bounded truth test for \(H\) is computable as well. Compute
\(f(e)\) and apply that test: the displayed equivalence would
decide \(K\), contrary to \cref{thm:halting-undecidable}.
Therefore HA does not prove (M), which is an actual instance of the
stated scheme with its decidability premise and universal parameter
closure intact.

The contradiction depends on total typed extracted functions and the
empty witness tuple of the double negation. Ordinary number realizability
uses a partial search program whose termination is promised only on
realizing inputs, so its validation of the same scheme is not in
conflict with this argument.
\end{proof}
\end{annexitem}
